\section{EDO de Flujo de Probabilidad}

\subsection{De EDP a EDO}

La EDP estocástica hacia adelante describe un proceso estocástico cuyas trayectorias son en zigzag (ya que las trayectorias del proceso de Wiener, aunque continuas, no son diferenciables en ninguna parte). Una pregunta natural es:

\textit{¿Existe una EDO determinista tal que la distribución marginal $p_t(\mathbf{x})$ que genera sea exactamente la misma que la de la EDP estocástica?}

\begin{importantbox}{EDO de Flujo de Probabilidad}
Dada la EDP estocástica hacia adelante:
$$
\mathrm{d}\mathbf{x} = f(\mathbf{x}, t)\,\mathrm{d}t + g(t)\,\mathrm{d}\mathbf{w}
$$
Existe una única EDO (llamada \textbf{EDO de Flujo de Probabilidad}):
$$
\mathrm{d}\mathbf{x} = \left[f(\mathbf{x}, t) - \frac{1}{2}g(t)^2 \nabla_{\mathbf{x}} \log p_t(\mathbf{x})\right]\mathrm{d}t
$$
De modo que la distribución marginal en cada instante $t$ de las trayectorias de esta EDO coincide con la distribución marginal $p_t(\mathbf{x})$ de la EDP estocástica.
\end{importantbox}

Observe la relación entre la EDO de Flujo de Probabilidad y la EDP estocástica hacia atrás:

\begin{table}[H]
\centering
\begin{tabular}{lcc}
\toprule
& \textbf{Corrección del término de arrastre} & \textbf{Término estocástico} \\
\midrule
EDP estocástica hacia atrás & $f - g^2 \nabla_{\mathbf{x}} \log p_t$ & $g(t)\,\mathrm{d}\bar{\mathbf{w}}$ \\
EDO de Flujo de Probabilidad & $f - \frac{1}{2} g^2 \nabla_{\mathbf{x}} \log p_t$ & Ninguno \\
\bottomrule
\end{tabular}
\caption{Comparación entre la EDP estocástica hacia atrás y la EDO de Flujo de Probabilidad}
\end{table}

En la forma de la EDO, el coeficiente del score cambia de $g^2$ a $\frac{1}{2}g^2$, y al mismo tiempo se \textbf{elimina el término del proceso de Wiener}.

\subsection{Ventajas de las EDOs}

La EDO de Flujo de Probabilidad ofrece varias ventajas prácticas:

\begin{enumerate}
    \item \textbf{Muestreo determinista}: Dado el ruido inicial $\mathbf{x}_T$, la solución de la EDO es única; el mismo ruido siempre genera la misma imagen. Esto es consistente con el comportamiento de DDIM (de hecho, DDIM puede verse como una discretización de la EDO de Flujo de Probabilidad).
    \item \textbf{Menos pasos de muestreo}: Las trayectorias de las EDOs deterministas son suaves, lo que permite utilizar tamaños de paso más grandes y solucionadores numéricos de orden superior (como el método de Runge-Kutta), reduciendo significativamente el número de pasos de muestreo.
    \item \textbf{Cálculo preciso de la verosimilitud}: Mediante la fórmula del cambio instantáneo de variables, es posible calcular con precisión el logaritmo de la verosimilitud $\log p(\mathbf{x}_0)$ de los datos utilizando la EDO.
    \item \textbf{Interpolación semántica}: Dado que la EDO establece una biyección entre el espacio de datos y el espacio de ruido, es posible realizar interpolaciones en el espacio latente para obtener resultados intermedios con significado semántico.
\end{enumerate}

\begin{knowledgebox}{La relación entre DDIM y la EDO de Flujo de Probabilidad}
DDIM puede entenderse como un esquema de discretización especial de la EDO de Flujo de Probabilidad correspondiente a VP-EDP. Esto explica dos características centrales de DDIM: (1) el proceso de muestreo es determinista (no hay término estocástico), y (2) se puede generar con alta calidad utilizando muchos menos pasos que durante el entrenamiento (por ejemplo, 50 pasos en lugar de 1000).
\end{knowledgebox}

\subsection{Trayectorias de EDO vs. Trayectorias de EDP}

\begin{itemize}
    \item \textbf{Trayectoria de EDP}: En zigzag, estocástica; diferentes trayectorias que parten del mismo punto llegan a destinos distintos. La distribución marginal se obtiene promediando sobre todas las trayectorias posibles.
    \item \textbf{Trayectoria de EDO}: Suave, determinista; cada punto inicial corresponde a un único punto final. La distribución marginal está dada por la deformación global del conjunto de trayectorias (flujo de probabilidad).
\end{itemize}

Ambas distribuciones marginales $p_t(\mathbf{x})$ son exactamente iguales en cada instante $t$, pero las \textbf{trayectorias individuales} son diferentes.

\begin{warningbox}{La EDO de Flujo de Probabilidad sigue requiriendo el Score}
Aunque la EDO elimina el término estocástico, la función score $\nabla_{\mathbf{x}} \log p_t(\mathbf{x})$ aparece aún en el término de arrastre. Por lo tanto, el proceso de entrenamiento no cambia: ya sea que se utilice un muestreador basado en EDP o uno basado en EDO para generar muestras, el modelo sigue entrenando la función score (o equivalentemente, predicción del ruido).
\end{warningbox}

\subsection{Relación con Neural ODE y Normalizing Flows continuos}

La EDO de Flujo de Probabilidad revela las conexiones profundas entre los modelos de difusión, las Neural ODE (Chen et al., 2018) y los Continuous Normalizing Flows (CNF).

\begin{knowledgebox}{Continuous Normalizing Flows}
Los CNF definen una EDO $\mathrm{d}\mathbf{x}/\mathrm{d}t = v_\theta(\mathbf{x}, t)$ mediante un campo vectorial parametrizado por una red neuronal, mapeando distribuciones simples (como la gaussiana) a distribuciones complejas de datos. La EDO de Flujo de Probabilidad es exactamente de esta forma: su campo vectorial es $f(\mathbf{x}, t) - \frac{1}{2}g(t)^2 \nabla_{\mathbf{x}} \log p_t(\mathbf{x})$, donde la parte del score se aproxima mediante una red neuronal.

La diferencia clave radica en que los CNF entrenan directamente el campo vectorial, mientras que los modelos de difusión primero entrenan una red score mediante score matching y luego utilizan la EDO de Flujo de Probabilidad para el muestreo. El entrenamiento del segundo enfoque es más estable porque el score matching tiene un objetivo de desruido explícito.
\end{knowledgebox}

\begin{importantbox}{Cambio Instantáneo de Variables}
Para un flujo generado por una EDO $\mathrm{d}\mathbf{x}/\mathrm{d}t = \mathbf{v}(\mathbf{x}, t)$, la evolución del logaritmo de la densidad satisface:
$$
\frac{\mathrm{d}\log p_t(\mathbf{x}_t)}{\mathrm{d}t} = -\mathrm{tr}\left(\frac{\partial \mathbf{v}}{\partial \mathbf{x}}\right)
$$
Sustituyendo el campo vectorial de la EDO de Flujo de Probabilidad en esta fórmula, es posible calcular con precisión el logaritmo de la verosimilitud $\log p_0(\mathbf{x}_0)$ de los datos. Esto representa una de las ventajas importantes de los modelos de difusión frente a las GAN: estas últimas no pueden proporcionar valores exactos de verosimilitud.
\end{importantbox}

\subsection{Resumen de la sección}

La EDO de Flujo de Probabilidad es el "espejo" determinista de la EDP estocástica: conserva la misma distribución marginal, pero con trayectorias suaves y sin aleatoriedad. Esto abre las puertas al muestreo determinista, a soluciones eficientes y al cálculo preciso de verosimilitudes, estableciendo además conexiones profundas entre los modelos de difusión, las Neural ODE y los Continuous Normalizing Flows.
